\section{The residual high--prime gate}
\label{sec:tpc44-next-gate}

The small--prime theorem and coefficient descent combine into a precise
sufficient condition.  Process every active prime \(p\le z_X\) first,
where \(z_X=(\log X)^{501/500}\), and call the resulting coefficient
energy \(E_{z_X}\).  If it is nonzero, continue through the high primes
and define
\begin{equation}
 \mathfrak A_{>z_X}
 :=\frac{E_R}{E_{z_X}}
 =\prod_{p_r>z_X}(1-\theta_r).
 \label{eq:tpc44-high-amplification}
\end{equation}
This is the exact high--face point--evaluation amplification.  It may
depend on the chosen order of the high primes through the individual
factors, but the ratio itself does not.

\begin{theorem}[A sufficient high--prime correlation gate]
\label{thm:tpc44-high-gate}
If the actual TPC equality coefficients satisfy
\begin{equation}
 \mathfrak A_{>z_X}\le X^{o(1)},
 \label{eq:tpc44-high-product-hypothesis}
\end{equation}
then
\begin{equation}
 Z_0(-\one)\le K_B\cD_0
 \label{eq:tpc44-conditional-endpoint}
\end{equation}
for all sufficiently large \(X\).  A sufficient local condition for
\eqref{eq:tpc44-high-product-hypothesis} is
\begin{equation}
 \sum_{p_r>z_X}(-\theta_r)_+=o(\log X).
 \label{eq:tpc44-negative-theta-sum}
\end{equation}
\end{theorem}

\begin{proof}
If \(E_{z_X}=0\), then \(F(-\one)=0\) and the result is immediate.
Otherwise, \cref{prop:tpc44-low-restriction-energy} and
\eqref{eq:tpc44-high-amplification} give
\begin{equation}
 |F(-\one)|
 \le\nu_{z_X}V^{1/2}\mathfrak A_{>z_X}^{1/2}.
 \label{eq:tpc44-final-F-bound}
\end{equation}
By \eqref{eq:tpc44-soft-variance},
\(
 V^{1/2}\le X^{o(1)}\cD_0
\), while
\cref{thm:tpc44-quantitative-freezing} gives
\begin{equation}
 \nu_{z_X}\le
 X^{1/400-12829/105210000+o(1)}.
 \label{eq:tpc44-nu-with-margin}
\end{equation}
The fixed power margin absorbs the factors \(X^{o(1)}\), as well as the
diagonal term \(\cD_0\), and proves the endpoint.

For \(-1\le\theta<1\),
\begin{equation}
 \log(1-\theta)\le(-\theta)_+.
 \label{eq:tpc44-log-theta-bound}
\end{equation}
If a factor has \(\theta=1\), the terminal energy vanishes.  Otherwise
summing \eqref{eq:tpc44-log-theta-bound} proves that
\eqref{eq:tpc44-negative-theta-sum} implies
\eqref{eq:tpc44-high-product-hypothesis}.
\end{proof}

The theorem is conditional, but the condition is now both exact and
coefficient specific.  It asks that high--prime restrictions not create
a polynomial cumulative anti--correlation.  A plausible route is to
prove a bilinear estimate for the paired coefficient sum
\begin{equation}
 \sum_m C_{r-1}(m)\overline{C_{r-1}(p_rm)}
 \label{eq:tpc44-paired-coefficient-target}
\end{equation}
uniformly after the preceding restrictions.  Merely proving that few
labels contain \(p_r\) gives, by Cauchy--Schwarz, a square--root density
loss; over many primes that is generally not summable.  The needed
input must also exploit an angle, sign, oscillation, or genuine
source--target bilinear structure.

The conditional face product gives an equivalent one--sided route.  If
\(A_z=0\), then repeated use of \eqref{eq:tpc44-face-doubling} gives
\(A_0=Z_0(-\one)=0\), so the endpoint is automatic.  If \(A_z>0\) and one
can prove for some cutoff \(z\)
\begin{equation}
 \log\frac{A_z}{\cD_0}
 +\sum_{p_r\le z}(\delta_r)_+
 \le\left(\frac1{400}+o(1)\right)\log X,
 \label{eq:tpc44-face-sufficient-condition}
\end{equation}
then \eqref{eq:tpc44-log-drift-criterion} closes the endpoint.  By
\eqref{eq:tpc44-Gram-face-increment}, each \(\delta_r\) is a normalized
real cross inner product between the \(p_r\)-free and
\(p_r\)-divisible blocks.  A bound of size
\(O(1/p_r)\) for its positive part would be summable with only a
polylogarithmic loss.  No such estimate is asserted here.

\subsection{What has and has not moved}

TPC--43 proved that almost every full prime--sign environment closes the
equality packet.  The present theorem proves more: after drawing only
the high signs, a density--one set of such environments permits arbitrary
adversarial overwriting of every prime up to
\((\log X)^{501/500}\).  This includes the physical values on a growing
block and is therefore a strict partial de--randomization.

The unresolved step is still substantial.  One must either prove
\eqref{eq:tpc44-high-product-hypothesis}, establish the face drift bound
\eqref{eq:tpc44-face-sufficient-condition}, derive a Type--II estimate
for \eqref{eq:tpc44-paired-coefficient-target}, or identify a fixed
structured class for \(C_X\) covered by a deterministic M\"obius
orthogonality theorem.  Failure of one route would be informative, but
none may be replaced by a suggestion that the all--minus high corner is
random.

\subsection{Claim ledger}

The unconditional results of this paper are:
\begin{enumerate}[label=(\roman*)]
 \item exact forward and reverse prime martingale layers;
 \item uniform compression of every small--prime face at cost \(\nu_z\);
 \item the Rankin bound and robust freezing through
 \((\log X)^{501/500}\);
 \item exact coefficient restriction and correlation products;
 \item exact conditional likelihood transport; and
 \item the vanishing--variance degree--two rough obstruction.
\end{enumerate}
The high--prime amplification bound is a clearly labeled sufficient
hypothesis, not a proved estimate for the physical coefficients.  No
result here proves the twin--prime conjecture, a prime--pair asymptotic,
fixed--shift four--M\"obius cancellation, or a general violation of sieve
parity.
